\chapter{Zipf's Law and Token Distributions}
\label{ch16}

Zipf's law is one of the oldest and most striking empirical regularities in language: a few token types are extremely common, while most are rare. This chapter does not pretend the mystery is solved. Its purpose is narrower:
\begin{itemize}
  \item review several major explanation families,
  \item keep careful track of what is empirical and what is proved,
  \item and explain what a future analytic-combinatorics program would actually need to model.
\end{itemize}

\section{Zipf's Law}

\begin{definition}
Take a corpus and rank token \emph{types} by decreasing empirical frequency. Zipf's law is the empirical claim that the relative frequency of rank $r$ is approximately
\[
  f(r) \approx \frac{C}{r^s},
\]
with $s$ often close to $1$.
\end{definition}

This is a statement about ranked \emph{types}. It is not yet a theorem, and it does not by itself say why such a pattern should arise.

A useful distinction is the following.
\begin{itemize}
  \item If the ranked frequencies satisfy $p_r \propto r^{-s}$, then the \emph{number of types} whose frequency exceeds a threshold $f$ is of order $f^{-1/s}$, obtained by solving $p_r \approx f$ for $r$.
  \item This is different from asking for the total probability mass carried by all types above that threshold. That is a token-weighted question rather than a type-counting one.
\end{itemize}
So one must be careful not to slide between type tails and token-mass tails.

\subsection{Mandelbrot's Refinement}

Mandelbrot proposed the shifted form
\[
  f(r) \approx \frac{C}{(r+b)^s},
\]
which often fits the highest-frequency ranks better. The shift mainly affects the head of the distribution rather than the long tail.

\section{A Useful Skeptic: Piantadosi}

Piantadosi's review \cite{Piantadosi2014} is valuable because it resists overinterpretation. Many very different models can be tuned to reproduce an approximate rank--frequency power law. So matching Zipf's law alone is weak evidence for any one mechanism.

That is the right methodological attitude for this chapter as well. A successful explanation should ideally connect Zipf behavior to other observables too: length, predictability, syntactic role, or temperature dependence.

\section{Pitman--Yor Processes}

A mathematically clean probabilistic explanation of heavy-tailed type frequencies comes from the Pitman--Yor process.

\subsection{Chinese Restaurant Process View}

With discount parameter $a \in (0,1)$ and concentration parameter $\theta > -a$, the Pitman--Yor process may be described by a Chinese restaurant process. After $n-1$ tokens have been generated, suppose type $i$ currently has count $c_i$, and $k$ distinct types have been seen so far. Then the next token is
\begin{itemize}
  \item an existing type $i$ with probability
  \[
    \frac{c_i-a}{n-1+\theta},
  \]
  \item or a brand-new type with probability
  \[
    \frac{\theta + ak}{n-1+\theta}.
  \]
\end{itemize}

The important qualitative facts are these.
\begin{itemize}
  \item The process produces heavy-tailed ranked frequencies.
  \item The vocabulary size grows polynomially:
  \[
    V(n) \asymp n^a
  \]
  in the standard asymptotic sense.
  \item The discount parameter $a$ controls how aggressively new types continue to appear.
\end{itemize}

For the present chapter, that last bullet is enough. We do not need a precise closed-form rank exponent formula to understand the role of the process. The main lesson is that a simple reinforcement-and-innovation mechanism can already produce long-tailed vocabularies.

\subsection{Hierarchical / Two-Stage Versions}

Hierarchical Pitman--Yor models, such as those studied by Goldwater, Griffiths, and Johnson \cite{GoldwaterGriffithsJohnson2011}, are more flexible. Linguistically, one may think of them as a cache or lexicon layer together with a token-generation layer. Those extra layers can fit natural-language data substantially better than a single flat process.

\section{Berman's Combinatorial View}

Berman's perspective is different in spirit. Rather than starting from a stochastic process prior, it starts from vocabulary growth. If the number of observed types after $n$ tokens behaves like
\[
  V(n) \sim n^\beta,
\]
then one may ask what rank--frequency laws are compatible with that growth.

The chapter does not need the full derivation here. What matters is the conceptual shift: the exponent is constrained by a combinatorial growth law, not introduced as a free fitting parameter. This makes Berman's story complementary to the Pitman--Yor one rather than obviously contradictory to it.

It is also worth being precise about what is already standard: polynomial vocabulary growth of the form $V(n) \asymp n^a$ is well known in Pitman--Yor settings. The genuinely interesting question is not whether that happens, but whether different explanatory frameworks are telling the same deeper story in different languages.

\section{Temperature Windows and Empirical Zipf Behavior}

Mikhaylovskiy \cite{Mikhaylovskiy2025b} studies how the empirical rank--frequency exponent changes with \emph{local decoding temperature}. This chapter emphasizes the adjective ``local'': by Chapter~15's distinctions, this is not the same as the global Gibbs temperature.

The empirical claim is that there is often a window of local temperatures in which generated text looks approximately Zipfian, while lower temperatures make the distribution too concentrated and higher temperatures make it too flat. That is an interesting empirical phenomenon, but it should not be read as a theorem about the global partition function without further argument.

\section{The Quantization Story}

Michaud et al. \cite{MichaudEtAl2023} propose a neural-mechanistic story: token usage may reflect a distribution over latent ``quanta'' or reusable features, and those latent usage frequencies can themselves be heavy-tailed. This is an appealing analogy to hierarchical Pitman--Yor models, but the connection is still heuristic rather than a formal equivalence theorem.

\section{A Simple Multinomial Proposition}

To see what Zipf's law does \emph{not} require, it is useful to record a very basic fact.

\begin{proposition}
Assume a finite vocabulary
\[
  \{w_1,\dots,w_V\}
\]
with strictly decreasing probabilities
\[
  p_1 > p_2 > \cdots > p_V > 0.
\]
If tokens are drawn i.i.d. from this distribution, then the empirical frequency of the type with rank $r$ converges almost surely to $p_r$.
\end{proposition}

\begin{proof}
By the strong law of large numbers, the empirical frequency of each type $w_i$ converges almost surely to $p_i$. Because the probabilities are strictly ordered, the empirical rank order stabilizes almost surely for large enough sample size, so the empirical frequency at rank $r$ converges to $p_r$.
\end{proof}

This proposition says that an i.i.d. model with Zipfian parameters produces Zipfian output, and one with non-Zipfian parameters does not. So the real explanatory question is not whether an i.i.d. model can display Zipf's law if hand-tuned that way; it is where the Zipf-like parameters come from in the first place.

\section{A Worked Mini-Example}

Consider a uniform unigram model on a vocabulary of size $V$. Then
\[
  p_r = \frac{1}{V}
  \qquad (1 \le r \le V),
\]
so the limiting rank--frequency curve is flat, not Zipfian. This reminds us that long tails do not appear automatically from simple independence; some additional structure is needed.

\section{A Speculative Bivariate Proposal}

The natural analytic-combinatorics temptation is to package token length and rank together in a two-variable series, for example
\[
  F(z,u) = \sum_{w \in \mathcal V} z^{|w|} u^{r(w)},
\]
where $r(w)$ is empirical rank.

This can be a useful \emph{formal summary}, but it is not a standard symbolic-method bivariate generating function. The reason is that rank is not a local structural parameter like ``number of leaves'' or ``number of occurrences of a letter.'' It is a global corpus statistic defined only after all types have been compared to each other.

So the honest status of such a two-variable series is:
\begin{itemize}
  \item interesting as a proposal,
  \item potentially useful in asymptotic models,
  \item but not something to which the earlier symbolic marking theorems apply directly.
\end{itemize}

If a future asymptotic model turned rank into a genuine tractable parameter, then singularity drift in the extra variable might encode a length--rank tradeoff. At present, that is a research direction rather than a theorem.

\section*{Summary}

This chapter's safest conclusions are:
\begin{enumerate}
  \item Zipf's law is empirical and does not by itself identify the right mechanism.
  \item Pitman--Yor, vocabulary-growth, and quantization stories are different explanation families, each illuminating a different aspect of the phenomenon.
  \item The strongest analytic-combinatorics proposal here is still speculative, because rank is not a standard local marked parameter.
\end{enumerate}

That is enough for the chapter's real role, which is to set up a sharper open problem rather than to pretend the theory is finished.

\section*{Further Reading}

\begin{itemize}
  \item \textbf{G. K. Zipf}, \emph{Human Behavior and the Principle of Least Effort} (1949) \cite{Zipf1949} --- the historical starting point.
  \item \textbf{S. T. Piantadosi}, ``Zipf's word frequency law in natural language: a critical review and future directions'' (2014) \cite{Piantadosi2014} --- the best skeptical overview.
  \item \textbf{S. Goldwater, T. L. Griffiths, and M. Johnson}, ``Producing power-law distributions and damping word frequencies with two-stage language models'' (2011) \cite{GoldwaterGriffithsJohnson2011} --- the standard hierarchical Pitman--Yor reference.
  \item \textbf{V. Berman}, ``Random text, Zipf's law, critical length, and implications for large language models'' (2025) \cite{Berman2025} --- for the vocabulary-growth viewpoint.
  \item \textbf{N. Mikhaylovskiy}, ``Zipf's law and the temperature window in transformer language models'' (2025) \cite{Mikhaylovskiy2025b} --- for the empirical temperature-window story.
  \item \textbf{E. J. Michaud et al.}, ``The quantization model of neural scaling laws'' (2023) \cite{MichaudEtAl2023} --- for the neural-mechanistic explanation family.
\end{itemize}

\section{Exercises}

\begin{exercise}
Give an example of a finite vocabulary distribution whose ranked probabilities are not close to a power law, and explain what its empirical rank--frequency plot converges to.
\end{exercise}

\begin{exercise}
In a finite vocabulary model with strictly decreasing probabilities $p_1>\cdots>p_V$, explain why the distinctness assumption is needed in the multinomial proposition.
\end{exercise}
